\section{Kategorie tensorowe, plecione, wstęgowe i modularne}\label{cat:sec:tensor-fusion-modular}

Rozwijamy pojęcia z rozdziałów~\ref{cat:sec:monoidal}, \ref{cat:sec:coherence} i~\ref{cat:sec:monoids-duality}, łącząc strukturę monoidalną z liniowością i strukturą abelową. Obiekty proste będą podstawowymi składnikami rozkładów iloczynów tensorowych. Dodatkowe struktury na dualnościach pozwolą definiować ślady również bez symetrii, co prowadzi do kategorii fuzyjnych, wstęgowych i modularnych.

\subsection{Sztywność i dokładność iloczynu tensorowego}

Zachowujemy konwencję z definicji~\ref{cat:def:dual-object}. Prawostronnią dualnością obiektu $X$ jest obiekt $X^*$ z koewaluacją i ewaluacją
\[
 i_X:I\to X\otimes X^*,\qquad e_X:X^*\otimes X\to I.
\]
Lewostronnią dualnością jest $L_X$ z morfizmami
\[
 \widetilde i_X:I\to L_X\otimes X,\qquad
 \widetilde e_X:X\otimes L_X\to I.
\]
Obie pary spełniają odpowiednie tożsamości trójkątne. Kategoria sztywna posiada takie dualności dla każdego obiektu (definicja~\ref{cat:def:rigid-compact}). Piszemy $I$ dla jednostki monoidalnej. W obliczeniach pomijamy izomorfizmy łączności i jednostki zgodnie z twierdzeniem o koherencji~\ref{cat:thm:mac-lane-coherence}; nie oznacza to dowolnego przestawiania czynników.

\begin{stwierdzenie}[Sprzężenia z obu stron]\label{cat:prop:tf-tensor-adjoints}
W kategorii sztywnej zachodzą sprzężenia
\[
 X^*\otimes-\dashv X\otimes-\dashv L_X\otimes-,
 \qquad
 -\otimes L_X\dashv-\otimes X\dashv-\otimes X^*.
\]
\end{stwierdzenie}
\begin{proof}
Pierwsze sprzężenie zostało wykazane w twierdzeniu~\ref{cat:thm:dual-adjunction}. Ponieważ $X$ jest prawostronnie dualny do $L_X$ dla pary $(\widetilde i_X,\widetilde e_X)$, to samo twierdzenie daje drugie sprzężenie w pierwszym ciągu. Dla sprzężeń w drugim ciągu przeprowadzamy rachunek po prawej stronie. Dla $f:A\otimes X\to B$ i $g:A\to B\otimes X^*$ definiujemy
\[
 \Phi(f)=(f\otimes\id_{X^*})(\id_A\otimes i_X),\qquad
 \Psi(g)=(\id_B\otimes e_X)(g\otimes\id_X).
\]
Z bifunktorialności i tożsamości trójkątnych otrzymujemy
\[
\begin{aligned}
 \Psi\Phi(f)
 &= f(\id_A\otimes[(\id_X\otimes e_X)(i_X\otimes\id_X)])=f,\\
 \Phi\Psi(g)
 &= (\id_B\otimes[(e_X\otimes\id_{X^*})(\id_{X^*}\otimes i_X)])g=g.
\end{aligned}
\]
Obie konstrukcje są naturalne w $A,B$, ponieważ składanie z morfizmem na wolnym czynniku komutuje z wstawianiem koewaluacji i stosowaniem ewaluacji. Daje to $-\otimes X\dashv-\otimes X^*$. Stosujemy ten wynik do $L_X$, którego prawostronnie dualnym obiektem jest $X$, otrzymując pozostałe sprzężenie.
\end{proof}

\begin{wniosek}[Dokładność iloczynu tensorowego]\label{cat:cor:tf-tensor-exact}
Jeżeli $\mathcal C$ jest abelową kategorią sztywną z addytywnym iloczynem tensorowym w każdej zmiennej, to $X\otimes-$ i $-\otimes X$ są funktorami dokładnymi.
\end{wniosek}
\begin{proof}
Oba funktory mają lewe i prawe sprzężone na mocy stwierdzenia~\ref{cat:prop:tf-tensor-adjoints}. Zachowują więc wszystkie istniejące granice i kogranice (twierdzenie~\ref{tw:adjoints-preserve}), w szczególności jądra i kojądra. Dokładność wynika z kryterium w twierdzeniu~\ref{cat:thm:exact-functor-kernels}.
\end{proof}

\begin{stwierdzenie}[Dualność jako antyrównoważność]\label{cat:prop:tf-duality-equivalence}
W kategorii sztywnej wybór prawych i lewych dualności wyznacza wzajemnie odwrotne z dokładnością do naturalnego izomorfizmu funktory kontrawariantne $X\mapsto X^*$ i $X\mapsto L_X$. Jeżeli kategoria jest abelowa, dualność przenosi ciągi dokładne na ciągi dokładne z odwróconym kierunkiem strzałek. Obiekty proste przechodzą na obiekty proste.
\end{stwierdzenie}
\begin{proof}
Funktorialność prawej dualności została udowodniona w stwierdzeniu~\ref{cat:prop:dual-functoriality}; dla lewej stosujemy ten sam dowód w kategorii z odwróconą kolejnością iloczynu tensorowego. Obiekt $X$ z parą $(i_X,e_X)$ jest lewostronnie dualny do $X^*$, a z parą $(\widetilde i_X,\widetilde e_X)$ jest prawostronnie dualny do $L_X$. Jednoznaczność dualności (stwierdzenie~\ref{cat:prop:dual-unique} i jego wersja lewostronna) daje kanoniczne izomorfizmy
\[
 L_{X^*}\cong X,\qquad (L_X)^*\cong X.
\]
Są naturalne: dla $f:X\to Y$ obie drogi w kwadracie naturalności spełniają tę samą równość z ewaluacją charakteryzującą morfizm dualny, jak w stwierdzeniu~\ref{cat:prop:dual-functoriality}. Jednoznaczność daje ich równość. Otrzymujemy antyrównoważność. Równoważność zachowuje granice i kogranice przez ich własności uniwersalne, a przejście do kategorii przeciwnej zamienia jądra z kojądrami. Stąd dokładność z odwróconymi strzałkami. Niezerowy obiekt prosty charakteryzuje się tym, że nie ma nietrywialnych podobiektów, równoważnie nietrywialnych obiektów ilorazowych. Antyrównoważność zamienia obie własności, więc zachowuje prostotę.
\end{proof}

\subsection{Skończone kategorie tensorowe i kategorie fuzyjne}

Od tej chwili ustalamy algebraicznie domknięte ciało $k$, chyba że wyraźnie zaznaczymy inaczej. Rozróżnienie między sztywnością, skończonością i półprostotą jest istotne: każda z tych własności odpowiada za inny aspekt teorii.

\begin{definicja}[Skończona kategoria abelowa nad ciałem]\label{cat:def:tf-finite-abelian}
Skończoną kategorią abelową nad $k$ nazywamy kategorię $k$-liniową równoważną $k$-liniowo kategorii skończenie wymiarowych lewych modułów nad pewną skończenie wymiarową asocjatywną algebrą z jedynką nad $k$. Kategoria $k$-liniowa ma przestrzenie morfizmów nad $k$ i dwuliniowe składanie. Nie wymagamy, aby równoważność z kategorią modułów zachowywała strukturę monoidalną.
\end{definicja}

\begin{stwierdzenie}[Konsekwencje skończoności]\label{cat:prop:tf-finite-abelian-properties}
W skończonej kategorii abelowej przestrzenie morfizmów są skończenie wymiarowe, każdy obiekt ma skończoną długość, istnieje skończenie wiele klas izomorfizmu obiektów prostych i dostatecznie wiele obiektów rzutowych. Dla obiektu prostego $S$ zachodzi $\operatorname{End}(S)\cong k$ jako algebra nad $k$.
\end{stwierdzenie}
\begin{proof}
W kategorii modułów nad skończenie wymiarową algebrą $A$ przestrzenie morfizmów są podprzestrzeniami przestrzeni odwzorowań liniowych. Każdy ścisły łańcuch podmodułów ma ograniczoną długość przez wymiar modułu, więc każdy niezerowy moduł ma ciąg kompozycyjny (wybieramy maksymalny podmoduł i stosujemy indukcję po wymiarze). Każdy moduł prosty $S$ jest ilorazem modułu regularnego $A$, ponieważ dla dowolnego $s\neq0$ odwzorowanie $a\mapsto as$ jest surjektywne. Ustalmy ciąg kompozycyjny $0=A_0\subsetneq\cdots\subsetneq A_n=A$. Dla surjekcji $q:A\to S$ wybieramy najmniejszy indeks $j$ taki, że $q(A_j)\neq0$. Wówczas indukowane odwzorowanie $A_j/A_{j-1}\to S$ jest niezerowym homomorfizmem między modułami prostymi, więc jest izomorfizmem. Wystarcza zatem skończona lista czynników jednego ciągu kompozycyjnego $A$. Każdy skończenie wymiarowy moduł jest ilorazem skończonej sumy kopii $A$, biorąc na przykład bazę liniową jako zbiór generatorów. Taki moduł wolny jest rzutowy: dla surjekcji i odwzorowania z modułu wolnego wybieramy podniesienia obrazów jego skończonej bazy modułowej. Daje to dostatecznie wiele obiektów rzutowych w sensie definicji~\ref{cat:def:enough-projectives-injectives}.

Każdy niezerowy endomorfizm $S$ jest odwracalny przez prostotę, zatem $D=\operatorname{End}(S)$ jest skończenie wymiarową algebrą z dzieleniem. Dla $d\in D$ liniowy operator lewego mnożenia przez $d$ na $D$ ma wartość własną $\lambda\in k$, ponieważ $k$ jest algebraicznie domknięte. Istnieje więc $v\neq0$ z $(d-\lambda\id_S)v=0$. Odwracalność $v$ daje $d=\lambda\id_S$. Stąd $D=k\id_S$. Wszystkie własności przenoszą się przez równoważność $k$-liniową.
\end{proof}

\begin{definicja}[Skończona kategoria tensorowa i kategoria fuzyjna]\label{cat:def:tf-finite-tensor-fusion}
Skończona kategoria tensorowa nad $k$ jest sztywną kategorią monoidalną, której kategoria bazowa jest skończoną kategorią abelową nad $k$, iloczyn tensorowy jest $k$-liniowy w każdej zmiennej, a jednostka $I$ jest prosta. W szczególności $\operatorname{End}(I)=k\id_I$. Kategorią fuzyjną nazywamy półprostą skończoną kategorię tensorową: każdy jej obiekt jest skończoną sumą prostą obiektów prostych. W części poświęconej kategoriom modularnym będziemy zwykle zakładać $\operatorname{char} k=0$; sama definicja użyta tutaj dopuszcza dowolną charakterystykę.
\end{definicja}

Nie wymagamy osobno dokładności iloczynu tensorowego, ponieważ wynika ona z wniosku~\ref{cat:cor:tf-tensor-exact}. Półprostota jest dodatkowym warunkiem, a nie konsekwencją sztywności.

\begin{stwierdzenie}[Nieznikanie iloczynu i obiekty rzutowe]\label{cat:prop:tf-projectives-tensor}
W skończonej kategorii tensorowej dla $X,Y\neq0$ zachodzi $X\otimes Y\neq0$. Ponadto iloczyn tensorowy obiektu rzutowego z dowolnym obiektem, w każdej kolejności, jest rzutowy.
\end{stwierdzenie}
\begin{proof}
Jeżeli $X\neq0$, to lewa koewaluacja $\widetilde i_X:I\to L_X\otimes X$ jest niezerowa: w przeciwnym razie tożsamość trójkątna dawałaby $\id_X=0$. Ponieważ $I$ jest prosty, $\widetilde i_X$ jest monomorfizmem. Dokładność $-\otimes Y$ daje monomorfizm
\[
 Y\cong I\otimes Y\hookrightarrow L_X\otimes X\otimes Y.
\]
Gdyby $X\otimes Y=0$, obiekt po prawej byłby zerowy przez addytywność iloczynu, co przeczy $Y\neq0$.

Niech $P$ będzie rzutowy. Na mocy stwierdzenia~\ref{cat:prop:tf-tensor-adjoints} mamy naturalne bijekcje
\[
 \Hom(P\otimes X,Z)\cong\Hom(P,Z\otimes X^*),\qquad
 \Hom(X\otimes P,Z)\cong\Hom(P,L_X\otimes Z).
\]
Dla epimorfizmu $Z\to W$ funktory po prawej stronie wysyłają go na epimorfizmy przez dokładność. Rzutowość $P$ zapewnia surjektywność indukowanych odwzorowań przestrzeni morfizmów. Korzystając z naturalności obu bijekcji i definicji obiektu rzutowego, dostajemy rzutowość $P\otimes X$ i $X\otimes P$.
\end{proof}

\begin{wniosek}[Kryterium półprostoty]\label{cat:cor:tf-unit-projective}
Skończona kategoria tensorowa jest fuzyjna wtedy i tylko wtedy, gdy jej jednostka monoidalna jest obiektem rzutowym.
\end{wniosek}
\begin{proof}
Jeśli jednostka $I$ jest rzutowa, to każdy obiekt $X\cong I\otimes X$ jest rzutowy przez stwierdzenie~\ref{cat:prop:tf-projectives-tensor}. Każdy krótki ciąg dokładny $0\to A\to B\to C\to0$ rozkłada się, ponieważ rzutowość $C$ daje prawą odwrotność epimorfizmu $B\to C$; stosujemy twierdzenie~\ref{cat:thm:abelian-split}. Indukcja po skończonej długości obiektu daje rozkład na sumę prostą obiektów prostych.

Odwrotnie w kategorii półprostej każdy epimorfizm $B\to C$ ma prawą odwrotność. Rozkładamy $B,C$ na skończone sumy proste obiektów prostych. Morfizmy między nieizomorficznymi obiektami prostymi są zerowe, a endomorfizmy prostych są skalarami przez stwierdzenie~\ref{cat:prop:tf-finite-abelian-properties}. Epimorfizm jest więc blokowym układem macierzy, osobno dla każdej klasy obiektów prostych. Każdy blok ma pełny rząd wierszowy, bo w przeciwnym razie istniałby niezerowy funkcjonał na odpowiednim bloku przeciwdziedziny znikający na obrazie, sprzecznie z epimorficznością. Wybieramy prawą odwrotność każdej takiej macierzy. Otrzymany morfizm jest prawą odwrotnością epimorfizmu. W szczególności każdy obiekt, w tym $I$, jest rzutowy.
\end{proof}

\subsection{Pierwsze przykłady}

\begin{przyklad}[Skończenie wymiarowe reprezentacje grupy skończonej]\label{cat:ex:tf-finite-group-representations}
Niech $G$ będzie grupą skończoną. Kategoria $\operatorname{Rep}_k(G)$ skończenie wymiarowych reprezentacji jest kategorią modułów nad algebrą grupową $k[G]$, zatem jest skończoną kategorią abelową. Iloczyn tensorowy ma działanie diagonalne $g(v\otimes w)=gv\otimes gw$, a jednostką jest jednowymiarowa reprezentacja trywialna. Dualność przestrzeni wektorowej z działaniem $(g\varphi)(v)=\varphi(g^{-1}v)$ jest dualnością reprezentacji: ewaluacja jest niezmiennicza względem działania $G$, a koewaluacja jest niezmiennicza, ponieważ pod identyfikacją $V\otimes V^*\cong\operatorname{End}_k(V)$ odpowiada identyczności, na której działanie przez sprzęganie jest trywialne. Tożsamości trójkątne są takie jak w twierdzeniu~\ref{cat:thm:vect-dualizable}. Symetria pochodzi z zamiany czynników. Otrzymujemy skończoną kategorię tensorową.

Gdy $\operatorname{char}k$ nie dzieli $|G|$, każdy podmoduł $W\subset V$ ma niezmiennicze dopełnienie. Wybieramy liniową projekcję $p:V\to W$ i definiujemy
\[
 \overline p(v)=\frac{1}{|G|}\sum_{g\in G}g\,p(g^{-1}v).
\]
Podstawienie $g=ha$ pokazuje $\overline p(hv)=h\overline p(v)$, a dla $w\in W$ każdy składnik jest równy $w$, zatem $\overline p|_W=\id_W$. Jądro $\overline p$ jest niezmienniczym dopełnieniem $W$. Indukcja po wymiarze daje półprostotę (twierdzenie Maschkego w tej sytuacji), więc kategoria jest fuzyjna.

Jeżeli natomiast $\operatorname{char}k=p>0$ i $p$ dzieli $|G|$, rozważmy epimorfizm reprezentacji $\varepsilon:k[G]\to k$, $\varepsilon(\sum_g a_g g)=\sum_g a_g$, gdzie $G$ działa na $k[G]$ przez lewe mnożenie, a na $k$ trywialnie. Niezmiennicze elementy $k[G]$ mają wszystkie współczynniki równe, czyli są wielokrotnościami $\sum_g g$. Ich obraz przez $\varepsilon$ jest zerowy, ponieważ $|G|=0$ w $k$. Nie istnieje więc prawa odwrotność $\varepsilon$ w kategorii reprezentacji. Jednostka nie jest rzutowa i z wniosku~\ref{cat:cor:tf-unit-projective} kategoria nie jest fuzyjna. Sztywność pozostaje jednak zachowana.
\end{przyklad}

\begin{przyklad}[Przestrzenie wektorowe z gradacją grupową]\label{cat:ex:tf-graded-vector-spaces}
Dla grupy skończonej $G$ rozpatrzmy kategorię $\operatorname{Vec}_G$ skończenie wymiarowych przestrzeni wektorowych $V=\bigoplus_{g\in G}V_g$ z morfizmami zachowującymi gradację. Jej iloczyn tensorowy definiujemy przez
\[
 (V\otimes W)_h=\bigoplus_{ab=h}V_a\otimes_k W_b.
\]
Jednostką jest $k$ skoncentrowane w stopniu neutralnym. Łączność mnożenia grupowego zapewnia, że zwykły izomorfizm łączności przestrzeni wektorowych zachowuje gradację: obie strony mają w stopniu $h$ składniki indeksowane trójkami $abc=h$. Izomorfizmy jednostki i diagramy koherencji są odziedziczone z przestrzeni wektorowych. Prawą i lewą dualność daje gradacja $(V^*)_g=(V_{g^{-1}})^*$: ewaluacja paruje wyłącznie stopnie wzajemnie odwrotne, a koewaluacja jest sumą koewaluacji poszczególnych składników. Obie mają stopień neutralny i spełniają tożsamości trójkątne stopień po stopniu.

Kategoria bazowa jest równoważna kategorii skończenie wymiarowych modułów nad produktem algebr $\prod_{g\in G}k$: ortogonalne idempotenty współrzędnych wyznaczają składniki gradacji i odwrotnie. Każdy obiekt rozkłada się na sumę prostą jednowymiarowych przestrzeni $k_g$ skoncentrowanych w jednym stopniu $g$. Są to dokładnie obiekty proste z dokładnością do izomorfizmu. Wobec tego $\operatorname{Vec}_G$ jest kategorią fuzyjną w każdej charakterystyce. Mamy
\[
 k_g\otimes k_h\cong k_{gh},\qquad (k_g)^*\cong k_{g^{-1}}.
\]
Gdy $G$ nie jest przemienna, nie istnieje plecenie: dla $g,h$ z $gh\neq hg$ obiekty $k_{gh}$ i $k_{hg}$ nie są izomorficzne, a plecenie musiałoby je utożsamiać. Gdy $G$ jest przemienna, zwykła zamiana czynników daje symetrię, ponieważ zachowuje gradację. Przykład ten oddziela fuzyjność od istnienia plecenia.
\end{przyklad}

\subsection{Pierścień Grothendiecka}

\begin{definicja}[Grupa i pierścień Grothendiecka]\label{cat:def:tf-grothendieck-ring}
Dla skończonej kategorii abelowej $\mathcal C$ grupa $K_0(\mathcal C)$ jest wolną grupą abelową na klasach izomorfizmu obiektów modulo relacje $[B]=[A]+[C]$ dla wszystkich krótkich ciągów dokładnych $0\to A\to B\to C\to0$. W przypadku skończonej kategorii tensorowej definiujemy dodatkowo
\[
 [X][Y]=[X\otimes Y],\qquad 1=[I].
\]
\end{definicja}
\begin{stwierdzenie}[Baza klas obiektów prostych]\label{cat:prop:tf-grothendieck-basis}
Niech $S_0=I,S_1,\ldots,S_{r-1}$ będą reprezentantami klas izomorfizmu obiektów prostych skończonej kategorii tensorowej. Wówczas $K_0(\mathcal C)$ jest wolną grupą abelową o bazie $[S_i]$. Mnożenie z definicji~\ref{cat:def:tf-grothendieck-ring} jest dobrze określone, łączne i ma jednostkę $[I]$. W szczególności
\[
 [X]=\sum_i m_i(X)[S_i],
\]
gdzie $m_i(X)$ jest krotnością czynnika prostego $S_i$ w ciągu kompozycyjnym $X$.
\end{stwierdzenie}
\begin{proof}
Istnienie ciągów kompozycyjnych wynika ze stwierdzenia~\ref{cat:prop:tf-finite-abelian-properties}. Dla kompletności uzasadnimy niezależność multizbioru czynników od ciągu. Wystarczy pracować w kategorii skończenie wymiarowych modułów nad algebrą z definicji~\ref{cat:def:tf-finite-abelian}, gdzie indukcję prowadzimy po wymiarze modułu. Dwa ciągi kompozycyjne kończą się maksymalnymi właściwymi podmodułami $U,V$ obiektu $X$. Jeśli $U=V$, stosujemy indukcję do $U$. W przeciwnym razie $U+V=X$ przez maksymalność, a dla $W=U\cap V$ mamy
\[
 U/W\cong X/V,\qquad V/W\cong X/U.
\]
Pierwszy izomorfizm otrzymujemy z surjekcji $U\to X/V$ o jądrze $W$, drugi analogicznie. Czynniki $U/W,V/W$ są proste. Wybieramy ciąg kompozycyjny $W$, przedłużamy go o $U$ lub $V$ i stosujemy indukcję do $U,V$. Obie listy czynników $X$ składają się z czynników $W$ oraz dwóch czynników $X/U,X/V$. To daje twierdzenie Jordana--Höldera w rozpatrywanym przypadku.

Dla krótkiego ciągu dokładnego ciąg kompozycyjny podobiektu można przedłużyć, biorąc przeciwobrazy kolejnych podmodułów ciągu kompozycyjnego ilorazu. Kolejne ilorazy tych przeciwobrazów są izomorficzne z odpowiednimi czynnikami ilorazu, gdyż indukowana surjekcja ma jądro będące poprzednim przeciwobrazem. Stąd $m_i(B)=m_i(A)+m_i(C)$. Funkcje $m_i$ definiują homomorfizm z $K_0(\mathcal C)$ do $\mathbb Z^r$. W drugą stronę wysyłamy standardową bazę na klasy $[S_i]$; relacje z kolejnych krótkich ciągów dokładnych ciągu kompozycyjnego dają wzór na $[X]$. Oba homomorfizmy są więc wzajemnie odwrotne.

Dokładność iloczynu tensorowego w obu zmiennych (wniosek~\ref{cat:cor:tf-tensor-exact}) zapewnia zgodność mnożenia z każdą relacją definiującą $K_0$. Łączność i jednostka wynikają z odpowiednich izomorfizmów monoidalnych. Nie potrzebujemy przemienności: kategoria może nie mieć plecenia.
\end{proof}

W kategorii niepółprostej pierścień Grothendiecka nie pamięta, czy obiekt jest sumą prostą swoich czynników kompozycyjnych. Rozróżnienie to jest istotne przy przejściu od skończonej kategorii tensorowej do kategorii fuzyjnej.

\subsection{Reguły fuzji i wzajemność Frobeniusa}

W tej podsekcji zakładamy półprostotę. Oznaczamy $b_i=[S_i]$ i przez $i^*$ indeks obiektu prostego izomorficznego z $S_i^*$; istnieje on na mocy stwierdzenia~\ref{cat:prop:tf-duality-equivalence}.

\begin{definicja}[Współczynniki fuzji]\label{cat:def:tf-fusion-coefficients}
Współczynniki $N_{ij}^{\ell}\in\mathbb Z_{\geq0}$ definiujemy przez
\[
 S_i\otimes S_j\cong\bigoplus_{\ell=0}^{r-1}
 S_\ell^{\oplus N_{ij}^{\ell}},\qquad
 b_i b_j=\sum_\ell N_{ij}^{\ell}b_\ell.
\]
\end{definicja}

\begin{stwierdzenie}[Opis współczynników przez przestrzenie morfizmów]\label{cat:prop:tf-fusion-hom}
Zachodzi
\[
 N_{ij}^{\ell}=\dim_k\Hom(S_i\otimes S_j,S_\ell)
 =\dim_k\Hom(S_\ell,S_i\otimes S_j).
\]
Ponadto $S_i^*\cong L_{S_i}$ i $(i^*)^*=i$.
\end{stwierdzenie}
\begin{proof}
W skończonej sumie prostej obiektów prostych każdy składnik izomorficzny z $S_\ell$ daje jednowymiarową przestrzeń morfizmów do lub z $S_\ell$, a pozostałe dają przestrzeń zerową. Wynika to ze stwierdzenia~\ref{cat:prop:tf-finite-abelian-properties} oraz faktu, że niezerowy morfizm między prostymi jest izomorfizmem. Otrzymujemy oba wzory.

Dla prostego $X$ i dowolnego prostego $Y$ sprzężenia ze stwierdzenia~\ref{cat:prop:tf-tensor-adjoints} dają
\[
 \Hom(I,X\otimes Y)\cong\Hom(X^*,Y),\qquad
 \Hom(X\otimes Y,I)\cong\Hom(Y,L_X).
\]
Przestrzenie po lewej mają ten sam wymiar przez już udowodniony opis krotności jednostki. Biorąc $Y=X^*$, dostajemy niezerowy morfizm $X^*\to L_X$ między prostymi, czyli izomorfizm. Z antyrównoważności wynika $L_{X^*}\cong X$, a więc $(X^*)^*\cong X$. Nie wybraliśmy jednak naturalnego monoidalnego izomorfizmu $X\to X^{**}$: jego istnienie na poziomie struktury monoidalnej jest odrębnym zagadnieniem.
\end{proof}

\begin{stwierdzenie}[Tożsamości fuzji]\label{cat:prop:tf-fusion-identities}
Dla wszystkich indeksów zachodzą
\[
\begin{aligned}
 N_{0j}^{\ell}&=N_{j0}^{\ell}=\delta_{j\ell},\\
 N_{ij}^{0}&=\delta_{j,i^*},\\
 \sum_a N_{ij}^a N_{ah}^{\ell}
 &=\sum_b N_{jh}^b N_{ib}^{\ell},\\
 N_{ij}^{\ell}&=N_{j^*i^*}^{\ell^*}.
\end{aligned}
\]
Ponadto zachodzi wzajemność Frobeniusa:
\[
 N_{ij}^{\ell}=N_{i^*\ell}^{j}=N_{\ell j^*}^{i}.
\]
\end{stwierdzenie}
\begin{proof}
Pierwsza równość wynika z izomorfizmów jednostki, druga ze sprzężeń użytych w dowodzie stwierdzenia~\ref{cat:prop:tf-fusion-hom}. Trzecia jest porównaniem współczynników przy $b_\ell$ w $(b_ib_j)b_h=b_i(b_jb_h)$. Czwarta wynika z kontrawariantnej dualności i izomorfizmu $(S_i\otimes S_j)^*\cong S_j^*\otimes S_i^*$ ze stwierdzenia~\ref{cat:prop:dual-tensor}. Dualność zachowuje skończone sumy proste jako antyrównoważność. Wreszcie
\[
 \Hom(S_i\otimes S_j,S_\ell)
 \cong\Hom(S_j,L_{S_i}\otimes S_\ell)
 \cong\Hom(S_j,S_{i^*}\otimes S_\ell),
\]
a sprzężenie po drugiej stronie daje
\[
 \Hom(S_i\otimes S_j,S_\ell)
 \cong\Hom(S_i,S_\ell\otimes S_j^*).
\]
Porównanie wymiarów przez stwierdzenie~\ref{cat:prop:tf-fusion-hom} daje obie równości. Wszystkie użyte bijekcje są liniowe, gdyż ich wzory są złożeniami morfizmów i iloczynów tensorowych z ustalonymi ewaluacjami i koewaluacjami.
\end{proof}

\begin{definicja}[Macierze fuzji]\label{cat:def:tf-fusion-matrices}
Na grupie z bazą $(b_0,\ldots,b_{r-1})$ oznaczamy przez $M_i$ macierz lewego mnożenia przez $b_i$, a przez $R_i$ macierz prawego mnożenia przez $b_i$. Używamy wektorów kolumnowych, zatem
\[
 (M_i)_{\ell j}=N_{ij}^{\ell},\qquad
 (R_i)_{\ell j}=N_{ji}^{\ell}.
\]
\end{definicja}
\begin{stwierdzenie}[Rachunek macierzy fuzji]\label{cat:prop:tf-fusion-matrices}
Macierze mają nieujemne całkowite współczynniki i spełniają
\[
 M_iM_j=\sum_\ell N_{ij}^{\ell}M_\ell,\qquad
 R_jR_i=\sum_\ell N_{ij}^{\ell}R_\ell,\qquad
 M_iR_j=R_jM_i.
\]
Ponadto
\[
 M_i^{\mathsf T}=M_{i^*},\qquad R_i^{\mathsf T}=R_{i^*}.
\]
Macierze $\sum_iM_i$ i $\sum_iR_i$ mają wszystkie współczynniki ściśle dodatnie.
\end{stwierdzenie}
\begin{proof}
Pierwsze trzy tożsamości są macierzowym zapisem łączności mnożenia. Transpozycje wynikają ze wzajemności Frobeniusa. Dla dowolnych $j,\ell$ obiekt $S_\ell\otimes S_j^*$ jest niezerowy na mocy stwierdzenia~\ref{cat:prop:tf-projectives-tensor}; ma więc pewien prosty składnik $S_i$. Wówczas $N_{\ell j^*}^{i}>0$, a z wzajemności $N_{ij}^{\ell}>0$. Daje to dodatniość współczynnika $(\ell,j)$ macierzy $\sum_i M_i$. Analogicznie wybierając składnik $S_i$ obiektu $S_j^*\otimes S_\ell$ dostajemy $N_{ji}^{\ell}=N_{j^*\ell}^{i}>0$, co dowodzi dodatniości sumy prawych macierzy.
\end{proof}

\subsection{Wymiar Frobeniusa--Perrona}

Poniższa konstrukcja jest rzeczywista i zależy wyłącznie od reguł fuzji, a nie od charakterystyki ciała $k$. Korzystamy z klasycznego twierdzenia Perrona--Frobeniusa w następującej postaci: macierz rzeczywista o wszystkich współczynnikach ściśle dodatnich ma dodatni promień spektralny, którego przestrzeń własna jest jednowymiarowa i zawiera wektor o wszystkich współrzędnych dodatnich. Por. N.~J.~Higham, \emph{What Is the Perron--Frobenius Theorem?}, 13 lipca 2021, twierdzenie dla macierzy nierozkładalnych. Nie zakładamy, że każda pojedyncza macierz fuzji jest nierozkładalna.

\begin{lemat}[Dodatni wektor własny macierzy nieujemnej]\label{cat:lem:tf-positive-eigenvector}
Jeśli $A$ ma współczynniki nieujemne i $Av=dv$ dla wektora $v$ o wszystkich współrzędnych dodatnich, to $d\geq0$ i $\rho(A)=d$, gdzie $\rho$ oznacza promień spektralny.
\end{lemat}
\begin{proof}
Z $Av\geq0$ wynika $d\geq0$. Dla dowolnej zespolonej wartości własnej $\lambda$ i odpowiadającego niezerowego wektora $w$ wybieramy indeks $p$ maksymalizujący $|w_p|/v_p$, oznaczając maksimum przez $t>0$. Mamy
\[
 |\lambda|t v_p
 =|(Aw)_p|
 \leq\sum_j A_{pj}|w_j|
 \leq t\sum_j A_{pj}v_j=tdv_p.
\]
Zatem $|\lambda|\leq d$; ponieważ $d$ jest wartością własną, dostajemy równość $\rho(A)=d$.
\end{proof}

\begin{twierdzenie}[Wymiar Frobeniusa--Perrona kategorii fuzyjnej]\label{cat:thm:tf-fpdim}
Istnieje dokładnie jeden homomorfizm pierścieni z jedynką
\[
 d:K_0(\mathcal C)\longrightarrow\mathbb R
\]
dodatni na klasach wszystkich obiektów prostych. Zachodzi $d(b_i)=\rho(M_i)$, $d(b_i)=d(b_{i^*})$ i $d(b_i)\geq1$. Dla obiektu $X$ definiujemy $\operatorname{FPdim}(X)=d([X])$.
\end{twierdzenie}
\begin{proof}
Niech $R=\sum_iR_i$. Macierz $R$ jest ściśle dodatnia przez stwierdzenie~\ref{cat:prop:tf-fusion-matrices}. Wybieramy dodatni wektor własny $v$ dla $\rho(R)$; jego przestrzeń własna jest jednowymiarowa. Ponieważ $M_iR=RM_i$, wektor $M_iv$ należy do tej przestrzeni własnej. Jest niezerowy: $M_i$ ma kolumnę odpowiadającą $b_0$, która jest wektorem jednostkowym o indeksie $i$, a $v_0>0$. Zatem $M_iv=d_iv$ z $d_i>0$ (współczynniki $M_i$ są nieujemne). Lemat~\ref{cat:lem:tf-positive-eigenvector} daje $d_i=\rho(M_i)$. Macierz $M_0$ jest identycznością, więc $d_0=1$, a stosując wzór mnożenia macierzy do $v$, dostajemy
\[
 d_id_j=\sum_\ell N_{ij}^{\ell}d_\ell.
\]
Przyporządkowanie $b_i\mapsto d_i$ rozszerza się więc do homomorfizmu pierścieni.

Jeśli $d'$ jest innym dodatnim homomorfizmem, to dla wektora kolumnowego $w$ o współrzędnych $w_i=d'(b_i)$ mamy
\[
 (M_i^{\mathsf T}w)_j=\sum_\ell N_{ij}^{\ell}w_\ell=d'(b_i)w_j.
\]
Lemat daje $d'(b_i)=\rho(M_i^{\mathsf T})=\rho(M_i)$, ponieważ macierz i jej transpozycja mają ten sam wielomian charakterystyczny. To dowodzi jednoznaczności. Równość $d_{i^*}=d_i$ wynika z $M_{i^*}=M_i^{\mathsf T}$. Wreszcie współczynnik $N_{ii^*}^0$ jest równy $1$, zatem
\[
 d_i^2=d_i d_{i^*}=1+\sum_{\ell\neq0}N_{ii^*}^{\ell}d_\ell\geq1.
\]
Dodatniość $d_i$ daje $d_i\geq1$.
\end{proof}

\begin{wniosek}[Własności wymiaru i obiekty odwracalne]\label{cat:cor:tf-fpdim-properties}
Wymiar Frobeniusa--Perrona jest addytywny względem sum prostych, multiplikatywny względem iloczynu tensorowego, zachowywany przez dualność i ściśle dodatni dla niezerowych obiektów. Każdy jego wymiar jest algebraiczną liczbą całkowitą. Obiekt prosty $S_i$ jest odwracalny względem iloczynu tensorowego wtedy i tylko wtedy, gdy $d_i=1$.
\end{wniosek}
\begin{proof}
Pierwsze własności wynikają z twierdzenia~\ref{cat:thm:tf-fpdim} i rozkładu $[X]$ w bazie. Dla dowolnego obiektu $X$ macierz lewego mnożenia przez $[X]$ ma współczynniki całkowite, a dodatni wektor $v$ z dowodu twierdzenia jest jej wektorem własnym z wartością $\operatorname{FPdim}(X)$. Wielomian charakterystyczny jest moniczny o współczynnikach całkowitych; otrzymujemy całkowitość algebraiczną (także dla $X=0$).

Jeśli $d_i=1$, równość z końca dowodu twierdzenia wymusza $N_{ii^*}^{\ell}=0$ dla $\ell\neq0$. Analogicznie dla $i^*$ mamy $N_{i^*i}^{\ell}=0$ dla $\ell\neq0$. Z półprostoty dostajemy $S_i\otimes S_i^*\cong I\cong S_i^*\otimes S_i$. Odwrotnie jeśli $S_i\otimes Y\cong I\cong Y\otimes S_i$, to $Y\neq0$ i $\operatorname{FPdim}(Y)\geq1$, gdyż każdy niezerowy obiekt ma co najmniej jeden prosty składnik o wymiarze co najmniej $1$. Multiplikatywność daje $1=d_i\operatorname{FPdim}(Y)$, więc $d_i=1$.
\end{proof}

\begin{definicja}[Wymiar Frobeniusa--Perrona kategorii]\label{cat:def:tf-fpdim-category}
Dla kategorii fuzyjnej definiujemy
\[
 \operatorname{FPdim}(\mathcal C)=\sum_{i=0}^{r-1}d_i^2.
\]
Jest to dodatnia rzeczywista algebraiczna liczba całkowita. Nie utożsamiamy jej bez dodatkowych założeń z wymiarem globalnym wyznaczonym przez ślady kategorialne.
\end{definicja}

\begin{przyklad}[Kategorie grupowe i reprezentacje]\label{cat:ex:tf-fpdim-group}
Dla $\operatorname{Vec}_G$ z przykładu~\ref{cat:ex:tf-graded-vector-spaces} macierze lewego mnożenia są macierzami permutacji elementów grupy. Wszystkie proste są odwracalne, więc $d_g=1$, $\operatorname{FPdim}(\operatorname{Vec}_G)=|G|$, a pierścień Grothendiecka jest pierścieniem grupowym $\mathbb Z[G]$.

Dla $\operatorname{Rep}_k(G)$ przy charakterystyce niedzielącej rzędu grupy zwykły wymiar liniowy jest dodatnim homomorfizmem pierścienia Grothendiecka do $\mathbb R$. Z jednoznaczności $\operatorname{FPdim}(V)=\dim_k V$. Aby obliczyć wymiar kategorii, rozkładamy lewy moduł regularny $k[G]$ na obiekty proste z krotnościami $a_i$. Przestrzeń $\Hom_{k[G]}(k[G],S_i)$ utożsamia się z $S_i$ przez ewaluację na jednostce (odwrotność przypisuje $s$ morfizm $a\mapsto as$). Z drugiej strony jej wymiar jest $a_i$ na mocy lematu Schura ze stwierdzenia~\ref{cat:prop:tf-finite-abelian-properties}. Mamy więc $a_i=\dim_k S_i$ i
\[
 |G|=\dim_k k[G]=\sum_i(\dim_k S_i)^2
 =\operatorname{FPdim}(\operatorname{Rep}_k(G)).
\]
\end{przyklad}

\begin{przyklad}[Reguła fuzji typu Fibonacciego]\label{cat:ex:tf-fibonacci-ring}
Rozważmy pierścień z bazą $1,x$ i relacją $x^2=1+x$, tj. $\mathbb Z[x]/(x^2-x-1)$ z wyróżnioną bazą klas $1,x$. Inwolucja dualności na bazie jest identycznością. Macierz lewego mnożenia przez $x$ w tej bazie i jej wielomian charakterystyczny są równe
\[
 M_x=\begin{pmatrix}0&1\\1&1\end{pmatrix},
 \qquad \det(tI_2-M_x)=t^2-t-1.
\]
Dodatni homomorfizm wymiarowy musi spełniać $d_x^2=1+d_x$, zatem
\[
 d_x=\frac{1+\sqrt5}{2},\qquad 1+d_x^2=\frac{5+\sqrt5}{2}.
\]
Pokazuje to, że formalne reguły fuzji dopuszczają niecałkowite wymiary, mimo całkowitości algebraicznej. Jeżeli dana kategoria fuzyjna realizuje ten pierścień z wyróżnioną bazą, powyższe obliczenie daje jej wymiary Frobeniusa--Perrona. Nie dowodzimy tu istnienia takiej kategoryfikacji ani nie utożsamiamy danych pierścienia z pełnymi danymi kategorii: izomorfizmy łączności i ich koherencje są dodatkową strukturą.
\end{przyklad}

\begin{uwaga}[Dalszy zakres]\label{cat:rem:tf-scope}
W następnych podsekcjach rozwiniemy plecenia, struktury piwotalne i sferyczne oraz ślady i skręcenia. Po tych przygotowaniach przejdziemy do centrum Drinfelda, centrum Mügera i kategorii modularnych.
\end{uwaga}
